\chapter{La Corruzione della Praxis}

\epigraph{``Come la citt\`a declina dall'aristocrazia alla timocrazia, dalla
timocrazia all'oligarchia, dall'oligarchia alla democrazia, dalla democrazia
alla tirannia, cos\`i fa l'anima.''}{\textcite{plato1997complete}, \emph{Repubblica} 543c--576b}

\section{Il Ciclo della Degenerazione}

Nella \emph{Repubblica}, Socrate traccia i regimi successivi con cui
una citt\`a giusta degenera in una ingiusta. Ogni regime rappresenta
una risposta diversa alla domanda ``Qual \`e il principio dominante
della citt\`a?'' L'onore (timocrazia), la ricchezza (oligarchia), la
libert\`a (democrazia) e la tirannia non sono meramente forme politiche
ma \emph{psichiche} --- ciascuna corrisponde a una diversa struttura
dell'anima, a una diversa risposta alla domanda ``Cosa governa l'essere
umano?''

Tracciamo la corrispondente degenerazione di \praxis. Quando l'azione
dimentica il suo fondamento in \physis, passa attraverso quattro regimi
sempre pi\`u degradati. Ciascuno \`e un modo di rispondere alla domanda
``Cosa governa l'azione?'' quando la natura non lo fa pi\`u.

\section{Timocrazia dell'Azione: Onore}

La prima corruzione \`e il regime dell'onore. Qui, l'azione \`e
governata dal desiderio di riconoscimento. Il praticante agisce non
perch\'e la natura lo richieda ma perch\'e altri vedranno e
approveranno. L'obiettivo non \`e la coerenza dell'azione ma la sua
visibilit\`a.

I segni del praticante timocratico: parla costantemente dei suoi
obiettivi ma raramente agisce su di essi. Ottimizza il suo albero di
obiettivi pubblico per l'impressione piuttosto che per la fattibilit\`a.
Colleziona strisce non per il ritmo che sostengono ma per il distintivo
che mostrano. Sceglie obiettivi per il loro prestigio --- ``Correr\`o
una maratona'' --- piuttosto che per la loro corrispondenza con la
natura reale del praticante.

Il regime timocratico \`e dove la piattaforma \praxis\ \`e pi\`u
facilmente pervertita. Le caratteristiche pubbliche --- strisce,
risultati, valutazioni dei partner --- possono diventare fini a s\'e
stesse, i burattini della caverna. La piattaforma deve essere progettata
per resistere a questo, per rendere la coerenza interna del praticante
pi\`u visibile a lui stesso della sua esposizione esterna.

\section{Oligarchia dell'Azione: Ricchezza}

La seconda corruzione \`e il regime della ricchezza. Qui, l'azione \`e
governata dal desiderio di guadagno materiale. Il praticante agisce solo
se l'azione produce un ritorno economico misurabile. Tutte le altre
azioni sono liquidate come inefficienti o irrilevanti.

I segni del praticante oligarchico: mantiene solo quegli obiettivi che
hanno un chiaro rendimento finanziario. Valuta ogni azione dal suo ROI.
Liquida i domini non economici dell'albero degli obiettivi --- spirito,
comunit\`a, gioco --- come indulgenze. Scambia la parte della natura
che produce ricchezza per l'intera natura.

Il regime oligarchico confronta la piattaforma \praxis\ con la sua sfida
pi\`u profonda. La piattaforma deve operare in un'economia di mercato;
deve far pagare i suoi servizi, e i suoi utenti la valuteranno
naturalmente in termini economici. La difesa della piattaforma \`e la
diversit\`a dei suoi domini di obiettivo: 14 domini di vita dal Corpo
e Fitness allo Spirito e Scopo, solo due dei quali (Sicurezza Finanziaria
e Ricchezza e Beni) sono economici. La stessa struttura dell'ontologia
afferma che l'azione umana ha pi\`u governatori del denaro.

\section{Democrazia dell'Azione: Libert\`a}

La terza corruzione \`e il regime della libert\`a. Qui, l'azione \`e
governata da nulla --- o piuttosto, dalla proposizione che tutti i
governatori sono ugualmente validi e nessuno ha autorit\`a. Il praticante
agisce su ogni impulso, persegue ogni desiderio, dichiara ogni obiettivo
--- e non ne completa nessuno.

I segni del praticante democratico: il suo albero di obiettivi \`e una
giungla di impegni concorrenti, nessuno dei quali persiste abbastanza a
lungo da dare frutti. \`E sedotto da ogni nuova possibilit\`a. Scambia
la proliferazione di obiettivi per la ricchezza della vita. Il suo
punteggio di coerenza oscilla selvaggiamente --- momenti di intenso
impegno seguiti da lunghi periodi di inerzia.

Il regime democratico \`e il pi\`u seducente perch\'e si presenta come
libert\`a. Il praticante crede di essere liberato dai vincoli della
natura, libero di scegliere qualsiasi azione in qualsiasi momento. Ma
questa \`e la libert\`a del centro commerciale --- la libert\`a di
scegliere tra merci, non la libert\`a di vivere secondo la propria
natura. Il praticante democratico non \`e libero; \`e disperso, diviso
contro s\'e stesso, l'Ego Frammentato della Trappola dello Specchio
\parencite{macintyre1981virtue}.

\section{Tirannia dell'Azione: Volont\`a}

La quarta corruzione \`e il regime della tirannia. Qui, l'azione \`e
governata dalla volont\`a nuda del praticante --- il desiderio di
imporre la propria visione sul mondo indipendentemente dai vincoli della
natura. Il praticante tirannico non chiede ``cosa richiede la natura?''
ma ``cosa voglio io?'' E vuole senza limite.

I segni del praticante tirannico: stabilisce obiettivi che violano i
limiti del proprio corpo, della propria comunit\`a, del proprio
ecosistema. Spinge oltre i segnali di fallimento come se il fallimento
fosse un nemico da sconfiggere piuttosto che un'informazione dalla
natura. Tratta il punteggio di coerenza del motore \physis\ come un
punteggio da massimizzare con ogni mezzo, inclusa la manipolazione del
punteggio stesso. La sua azione diventa compulsiva, guidata, distruttiva
--- sia di s\'e che del mondo intorno a lui.

Il regime tirannico \`e il punto finale della degenerazione. Qui,
\praxis\ ha completamente dimenticato il suo fondamento in \physis. Lo
strumento \`e diventato il padrone. Il praticante crede di essere pi\`u
potente quando \`e pi\`u autonomo dalla natura; in realt\`a, \`e pi\`u
schiavo, perch\'e la sua volont\`a non \`e guidata da nessuna
intelligenza oltre s\'e stessa.

\section{Il Meccanismo del Declino}

Come declina \praxis\ attraverso questi regimi? Non per una forza esterna
ma per una logica interna. Ogni regime contiene il seme del successivo.

Il praticante che agisce per onore (timocrazia) scopre che l'onore non
\`e un bene durevole --- dipende dalle opinioni degli altri, che sono
volubili. Cos\`i si rivolge a qualcosa di pi\`u solido: la ricchezza
(oligarchia). Ma la ricchezza, una volta accumulata, non soddisfa il
bisogno di significato --- crea solo pi\`u desiderio. Cos\`i si rivolge
alla libert\`a (democrazia), perseguendo ogni desiderio indiscriminatamente.
Ma questo lo disperde, e la dispersione crea un bisogno disperato di
unit\`a. In questo vuoto entra il tiranno --- la volont\`a che unifica
dominando tutto sul suo cammino.

Il meccanismo \`e lo stesso in ogni stadio: \praxis\ cerca un governatore,
e ogni falso governatore fallisce. L'onore fallisce, la ricchezza
fallisce, la libert\`a fallisce --- solo la volont\`a del tiranno sembra
avere successo, e anche quel successo \`e il pi\`u grande fallimento di
tutti, perch\'e il tiranno \`e governato dalla propria volont\`a, che
non ha limite e non pu\`o mai essere soddisfatta.

\section{L'Antidoto}

L'antidoto a questo declino non \`e un quinto regime migliore. \`E il
riconoscimento che la ricerca di un governatore umano dell'azione ---
onore, ricchezza, libert\`o o volont\`a --- \`e essa stessa l'errore.
L'unico governatore adeguato dell'azione \`e la natura stessa (\physis),
non perch\'e la natura sia pi\`u forte della volont\`a del tiranno (lo
\`e), ma perch\'e la natura \`e la fonte dell'intelligibilit\`a
dell'azione. Solo quando l'azione \`e governata dalla natura ha un fine
oltre s\'e stessa, un criterio di successo che non \`e arbitrario, un
fondamento che non cambia con ogni mutamento di moda o fortuna.

La piattaforma \praxis\ e il motore \physis\ sono costruiti su questo
riconoscimento. Non creano il governo dell'azione da parte della natura;
lo rivelano. Il punteggio di coerenza non \`e un'invenzione umana ma
una scoperta naturale --- la traduzione del giudizio stesso della natura
in una forma che i praticanti possano vedere e da cui imparare.

Il declino in timocrazia, oligarchia, democrazia e tirannia non \`e
inevitabile. \`E una conseguenza del dimenticare che \praxis\ \`e lo
strumento di \physis. Il praticante che lo ricorda --- che coltiva le
discipline, presta attenzione alla coerenza e torna a testimoniare per
gli altri --- ha rotto il ciclo.

\begin{figure}[ht]
\centering
\begin{tikzpicture}[
  reg/.style={draw, rounded corners, align=center, minimum width=28mm,
              minimum height=8mm, font=\small},
  >={Stealth[length=2mm]}]
  \node[reg, fill=green!15] (praxis) {Praxis governata da Physis};

  \node[reg, fill=yellow!15, below=of praxis, yshift=-4mm] (timo)
    {Timocrazia\\\tiny governata dall'Onore};
  \draw[->, thick] (praxis) -- node[right, font=\tiny] {oblio} (timo);

  \node[reg, fill=orange!15, below=of timo, yshift=-4mm] (olig)
    {Oligarchia\\\tiny governata dalla Ricchezza};
  \draw[->, thick] (timo) -- node[right, font=\tiny] {l'onore fallisce} (olig);

  \node[reg, fill=red!15, below=of olig, yshift=-4mm] (demo)
    {Democrazia\\\tiny governata dalla Libert\`a};
  \draw[->, thick] (olig) -- node[right, font=\tiny] {la ricchezza fallisce} (demo);

  \node[reg, fill=red!30, below=of demo, yshift=-4mm] (tyr)
    {Tirannia\\\tiny governata dalla Volont\`a};
  \draw[->, thick] (demo) -- node[right, font=\tiny] {la libert\`a disperde} (tyr);

  \draw[->, thick, cohgreen, dashed] (tyr.east) -- ++(20mm,0)
    node[right, font=\tiny, cohgreen] {Antidoto: ritorno a Physis};
\end{tikzpicture}
\caption{La corruzione quadruplice di \praxis. Ogni regime \`e governato da un falso
principio che fallisce, aprendo la strada al successivo. L'antidoto \`e il ritorno
alla natura come unico governatore adeguato.\label{fig:corruzioni}}
\end{figure}
